\section*{Orientaciones bibliográficas}
\addcontentsline{toc}{section}{Orientaciones bibliográficas}
Para anillos y homomorfismos, véase \citet[cap.~III, 1.1/2/5/8--15/19 y 2.1--3]{aluffi2009}. Los polinomios y la evaluación se tratan en III.2.1/4--5.

Para productos, ideales, cocientes y polinomios, véase \citet[cap.~III, 2.4--7, 3.1/3/8--11 y 4.1--6]{aluffi2009}. Los ejercicios corresponden a III.2.13, III.3.2 y III.4.1.

Para ideales primos y maximales, véase \citet[cap.~III, 4.6--7/9--14]{aluffi2009}; la definición de álgebra se encuentra en III.5.1--7. Los ejercicios sobre $K[x]$ y $\mathbb C[x]$ son III.4.4 y III.4.21; para el segundo se usa el teorema fundamental del álgebra.

Para factorización y dominios principales o euclidianos, véase \citet[cap.~V, 1.7--12 y 2.1--8]{aluffi2009}. El ejemplo de infinitas variables corresponde al ejercicio V.1.15; el caso de polinomios sobre un campo retoma III.4.4.

Para Euclides y la aplicación de Zorn a ideales maximales, véase \citet[cap.~V, 2.9--12 y 3.1/3/5]{aluffi2009}; para contenido y primitividad, V.4.1--9. Las equivalencias fundacionales se enuncian sin prueba. Los ejercicios corresponden a V.3.9 y V.4.1.

Para fracciones y factorialidad de polinomios, véase \citet[cap.~V, 4.10--18]{aluffi2009}. Los ejercicios de elección de factores e infinitas variables se relacionan con V.4.23 y V.1.15.

Para raíces, extensiones e irreducibilidad, véase \citet[cap.~V, apartados 5.1--4 y resultados 5.1--17]{aluffi2009}. El teorema fundamental del álgebra se usa sin prueba. Los ejercicios de grado impar y campos finitos corresponden a V.5.17 y V.5.11, respectivamente; la proposición V.5.17 es el criterio de Eisenstein.

La propiedad universal de las álgebras conmutativas libres se desarrolla en \citet[cap.~III, 5.18 y apartados 6.1--3]{aluffi2009}. Para divisibilidad e irreducibilidad, véanse V.1.1--6 y el ejercicio V.1.12. El teorema chino del resto se aplica en VI.5.1--7.
